Os \emph{sistemas de prova} fornecem métodos puramente sintáticos para
caracterizar consequência e compatibilidade entre sentenças.
O \emph{cálculo de sequentes} é um desses sistemas de
prova. Uma \emph{derivação} nele consiste em uma árvore de sequentes (um
sequente $\Gamma \Sequent \Delta$ consiste em duas sequências de fórmulas
separadas por~$\Sequent$). Os sequentes mais ao topo em uma derivação
são \emph{sequentes iniciais} da forma $!A \Sequent !A$.  Todos os demais
sequentes, para que a derivação seja correta, devem ser corretamente
justificados por uma de várias \emph{regras de inferência}. Elas vêm em
pares; uma regra para operar no lado esquerdo e no lado direito de um
sequente para cada conectivo e quantificador. Por exemplo, se um sequente
$\Gamma \Sequent \Delta, !A \lif !B$ é justificado pela regra
$\RightR{\lif}$, o sequente precedente (a \emph{premissa}) deve
ser $!A, \Gamma \Sequent \Delta, !B$. Algumas regras também permitem que a ordem
ou o número de sentenças em um sequente seja manipulado, por exemplo,
a regra \RightR{\Exchange} permite que duas fórmulas do lado direito de um
sequente sejam trocadas.

Se há uma derivação do sequente $\quad \Sequent !A$, dizemos que
$!A$ é um \emph{teorema} e escrevemos~$\Proves !A$. Se há uma
derivação de $\Gamma_0 \Sequent !A$ em que todo $!B$ em $\Gamma_0$ está
em~$\Gamma$, dizemos que $!A$ é \emph{derivável de}~$\Gamma$ e escrevemos
$\Gamma \Proves !A$. Se há uma derivação de
$\Gamma_0 \Sequent \quad$ em que todo $!B$ em $\Gamma_0$ está
em~$\Gamma$, dizemos que $\Gamma$ é \emph{inconsistente}; caso contrário,
\emph{consistente}. Essas noções estão inter-relacionadas; por exemplo, $\Gamma
\Proves !A$ sse $\Gamma \cup \{\lnot !A\}$ é inconsistente. Elas estão
também relacionadas às noções semânticas correspondentes; por exemplo, se $\Gamma
\Proves !A$, então $\Gamma \Entails !A$. Essa propriedade dos sistemas de
prova---o que pode ser derivado de $\Gamma$ tem garantia de ser
consequência de~$\Gamma$---é chamada de \emph{correção}. O \emph{teorema da
correção} é provado por indução sobre o comprimento das derivações, mostrando
que cada inferência individual preserva a validade do sequente
conclusão, desde que os sequentes premissa sejam válidos.
